\documentclass[11pt]{article}

\usepackage[margin=1in]{geometry}
\usepackage{amsmath,amssymb,amsthm}
\usepackage{graphicx}
\usepackage[hidelinks]{hyperref}

\newtheorem{theorem}{Theorem}
\newtheorem{lemma}[theorem]{Lemma}

\newcommand{\R}{\mathbb R}
\newcommand{\Pbb}{\mathbb P}
\newcommand{\E}{\mathbb E}
\newcommand{\eps}{\varepsilon}
\newcommand{\norm}[1]{\left\lVert #1 \right\rVert}

\title{A Cauchy-Stable Partial Separation Reduction for the \(B_1^n\) Case}
\author{arXiv automatic math run \texttt{fa\_banach\_001}}
\date{July 3, 2026}

\begin{document}
\maketitle

\section*{Status}

This packet gives a candidate partial result, likely valid pending expert
review, for the Part 1 separation dimension-reduction challenge in
Mendelson--Milman--Paouris \cite{MMP}.  It does not complete the authors'
program: the separation loss obtained below is \(O(k\log k)\), whereas the
program needs a universal constant loss.

\section*{Source Challenge}

The source paper is Shahar Mendelson, Emanuel Milman, and Grigoris Paouris,
\emph{Generalized Dual Sudakov Minoration via Dimension Reduction - A Program},
arXiv:1610.09287.  In the conclusion, after noting that Part 2 is essentially
settled up to the slicing problem, the authors write that the main remaining
challenge is Part 1, and specifically mention \(K=B_1^n\).

\begin{center}
\includegraphics[width=\textwidth]{figures/open_problem_crop.png}
\end{center}

Part 1' asks, informally, for a linear map to dimension \(m\simeq k\) that
keeps \(e^k\) many \(K\)-separated points separated while sending a positive
amount of the original measure into a sufficiently massive target body.  The
authors suggest that for bodies such as \(B_1^n\), one may need random maps
distributed differently from Haar-random Euclidean projections.

\section*{Proof Intuition}

For \(K=B_1^n\), the natural non-Haar random map is a Cauchy matrix.  The
standard Cauchy law is \(1\)-stable: every row applied to a vector \(x\) has
scale \(\norm{x}_1\).  Thus a difference with \(\ell_1\)-norm \(>1\) has an
exponentially small probability of landing in a fixed small \(\ell_1^m\)-ball,
so a union bound preserves \(e^k\) separations in \(m\simeq k\) dimensions.
On the other hand, vectors in \(B_1^n\) are mapped into an
\(O(m\log m)B_1^m\) ball with constant probability, giving the required
massiveness after averaging over the measure.  The \(m\log m\) radius is the
price paid for the heavy Cauchy tails, and is exactly why this remains only a
partial result.

\section*{The Partial Result}

\begin{theorem}[Cauchy one-sided reduction for \(B_1^n\)]
\label{thm:cauchy-b1}
There are universal constants \(a,A,c>0\) such that the following holds.  Let
\(k\ge 1\), \(M\le e^k\), and let \(x_1,\ldots,x_M\in\R^n\) satisfy
\[
  \norm{x_i-x_j}_1>1 \qquad (i\ne j).
\]
Let \(\mu\) be a Borel probability measure on \(\R^n\) with
\(\mu(B_1^n)\ge e^{-q}\), where \(0\le q\le k/2\).  Put
\[
  m=\lceil k\rceil, \qquad L=A\,m\log(e m)\,B_1^m .
\]
Then there is a linear map \(T:\R^n\to\R^m\) such that
\[
  \norm{Tx_i-Tx_j}_1>a \qquad (i\ne j)
\]
and
\[
  T_*\mu(L)\ge c e^{-q}.
\]
Equivalently, the images are \((1/R)L\)-separated with
\[
  R\le C\,m\log(e m).
\]
\end{theorem}

\begin{lemma}[Two elementary Cauchy estimates]
\label{lem:cauchy}
Let \(G_1,\ldots,G_m\) be independent standard Cauchy random variables.  There
are universal constants \(\eta,A>0\) such that, for every \(m\ge 1\),
\[
  \Pbb\left\{\sum_{j=1}^m |G_j|\le \eta\right\}\le e^{-6m}
\]
and
\[
  \Pbb\left\{\sum_{j=1}^m |G_j|\le A m\log(e m)\right\}\ge \frac12 .
\]
\end{lemma}

\begin{proof}
For \(0<u\le 1\), \(\Pbb\{|G_1|\le u\}\le C u\).  Choosing \(\eta>0\) so that
\(C\eta\le e^{-6}\), the first estimate follows from the inclusion
\(\{\sum |G_j|\le\eta\}\subseteq\{|G_j|\le\eta \text{ for all }j\}\).

For the second estimate, use
\[
  \Pbb\{|G_1|>u\}\le \frac{C}{u},
  \qquad
  \E\bigl[|G_1|\mathbf 1_{\{|G_1|\le u\}}\bigr]\le C\log(eu).
\]
With \(R=A m\log(e m)\),
\[
\begin{aligned}
\Pbb\left\{\sum |G_j|>R\right\}
&\le \Pbb\left\{\max_j |G_j|>R\right\}
   +\Pbb\left\{\sum |G_j|\mathbf 1_{\{|G_j|\le R\}}>R\right\}  \\
&\le \frac{C m}{R}+\frac{C m\log(eR)}{R}.
\end{aligned}
\]
Taking \(A\) large enough makes the last quantity at most \(1/2\), uniformly
in \(m\).
\end{proof}

\begin{proof}[Proof of Theorem~\ref{thm:cauchy-b1}]
Let \(T=(g_{jr})_{1\le j\le m,1\le r\le n}\), where the \(g_{jr}\) are
independent standard Cauchy random variables.  The characteristic function
identity \(\E e^{itG}=e^{-|t|}\) implies the \(1\)-stability relation
\[
  (Tx)_j \stackrel{d}{=} \norm{x}_1 G_j ,
\]
with independent \(G_j\)'s as \(j\) varies.

If \(h\notin B_1^n\), then \(\norm{h}_1>1\), and Lemma~\ref{lem:cauchy} gives
\[
  \Pbb\{\norm{Th}_1\le \eta\}
  \le
  \Pbb\left\{\sum_{j=1}^m |G_j|\le \eta\right\}
  \le e^{-6m}.
\]
Therefore,
\[
  \Pbb\{\exists i<j:\norm{T(x_i-x_j)}_1\le \eta\}
  \le M^2 e^{-6m}\le e^{2k-6m}\le e^{-4k}.
\]

For \(x\in B_1^n\), the same stability relation and the second part of
Lemma~\ref{lem:cauchy} imply
\[
  \Pbb\{Tx\in A m\log(e m)B_1^m\}\ge \frac12 .
\]
Consequently, for \(L=A m\log(e m)B_1^m\),
\[
  \E\bigl[T_*\mu(L)\bigr]
  =\int \Pbb\{Tx\in L\}\,d\mu(x)
  \ge \frac12 \mu(B_1^n)
  \ge \frac12 e^{-q}.
\]
Since \(0\le T_*\mu(L)\le 1\), this implies
\[
  \Pbb\{T_*\mu(L)\ge \tfrac14 e^{-q}\}\ge \tfrac14 e^{-q}.
\]
Because \(q\le k/2\), the latter probability is larger than \(e^{-4k}\) after
adjusting constants.  Hence there is a realization of \(T\) for which both the
separation event and the massiveness event occur.  Taking \(a=\eta\) and
\(c=1/4\) proves the theorem.  Finally,
\[
  \eta B_1^m=\frac{\eta}{A m\log(e m)}L,
\]
so the separation may be written as \((1/R)L\)-separation with
\(R\le C m\log(e m)\).
\end{proof}

\section*{Limitations and Novelty Check}

The result is not the desired Part 1' for \(B_1^n\), because \(R\) grows like
\(k\log k\).  The source program needs \(R\) independent of \(k,n\) (or at
least a substantially milder loss that can be absorbed in the bootstrap).  The
packet should therefore be read as evidence that \(1\)-stable, non-Haar random
maps are relevant to the \(B_1^n\) challenge, not as a completion of the
program.

Novelty check performed on July 3, 2026: local run indexes were searched for
arXiv:1610.09287, ``Massive Separation Dimension Reduction'', \(B_1^n\), Cauchy,
and one-sided Johnson--Lindenstrauss terms; no duplicate packet was found.
Web searches for ``one-sided Johnson-Lindenstrauss Cauchy l1'', ``Cauchy
dimension reduction l1 separation'', and related stable-projection phrases did
not surface a direct match.  The result uses only elementary Cauchy stability
and is not claimed to be absent from the broader algorithmic dimension
reduction literature.

\section*{Human Review Recommendation}

Reviewers should check whether the \(O(k\log k)\) loss can be reduced by using
a different target body, truncating the Cauchy map, or replacing the \(\ell_1^m\)
massiveness body by a Lorentz or weak-\(\ell_1\) body compatible with Part 2.
As written, the proof is self-contained and the main risk is relevance rather
than correctness.

\begin{thebibliography}{9}

\bibitem{MMP}
S. Mendelson, E. Milman, and G. Paouris,
\emph{Generalized Dual Sudakov Minoration via Dimension Reduction - A Program},
arXiv:1610.09287, 2016.

\end{thebibliography}

\end{document}
